\chapter{A quantum ontology}
In this chapter, we present a complete example of rewriting, starting from an ontology expressed in OWL (Section~\ref{sec:owl}), and arriving at a Prolog formulation suitable for the pipeline presented in the previous chapter.

Following the literature, we begin with the translation of the ontology into Prolog, discussing how to translate concepts such as classes, subclasses, individuals, and relations. We then focus on the incompatibility of this first version with QA-Prolog, highlight the reasons for this incompatibility, and propose possible solutions.

Finally, we present some examples to demonstrate the equivalence of the two transformations and provide observations on the final result obtained after the transformations performed by QA-Prolog.

\section{Starting ontology}
The ontology we will translate in Prolog is not the one presented in Chapter~\ref{ch:ont} but is a more complex in order to show an example more similar to a real ontology. 

The domain of interest of our ontology is the animal kingdom; this domain serves as a pretext for a deeper description than the one shown in the family tree example (Section~\ref{sec:ont_ex}). We will see how to describe animal species, their feeding habits, and the relationships between individuals.

\subsection{Ontology structure}
In this ontology, we describe a small part of the animal kingdom. The more complex class hierarchy represents the taxonomy of animals; it characterizes living beings with increasing precision until identifying their species.

The other classes instead focus feeding habits of animals. We distinguish between prey and predators, and between carnivorous and herbivorous animals (in this small example, for simplicity, we ignore omnivores\footnote{The reason will be clear in the next section.}).

Finally, we consider relationships between animals. We define two types of relationships: parental relation that connect two animals in parent–son roles, and relation related to feeding habits that link two animals in the roles of predator and prey.

The T-Box of our ontology can be seen in Figure~\ref{fig:ont_animal}, we define also some individual and their relation to populate our ontology\footnote{The code for this ontology is available at \url{URL}}.
\begin{figure}[h]
	\centering
	\includegraphics[width=.9\linewidth]{setup/overview-image}
	\caption{Ontology T-Box}
	\label{fig:ont_animal}
\end{figure}

\subsection{Ontology constraints}
In this section, we present the constraints on classes and individuals represented in our ontology. These constraints allow us to make inference on the ontology and derive new information from what is already present, also thanks to these constraints we can determine if all the information contained in the ontology is consistent with the theoretical model we are describing.

Constrains defined in animal ontology are:
\begin{itemize}
	\item Constrain on classes:
	\begin{itemize}
		\item different spices are disjoint classes;
		\item  carnivore and herbivore are disjoint classes;
	\end{itemize}
	\item Constrain on individual:
	\begin{itemize}
		\item If animal $A$ hunt $B$, $A$ is a predator and carnivore and $B$ is a prey.
	\end{itemize}
\end{itemize}

From defined constraints we can see that a prey and predator are not disjoint classes and be a prey does not automatically mean that the animal is a herbivore. This is coherent with the animal kingdom where a certain animal could be both a prey and a predator.

It should also be clear why we have omitted the omnivores from our description, this simplification allows us to describe herbivore and carnivore as disjoint and to infer that an animal that hunt is automatically a carnivore.

\subsection{Ontology contents}
We now examine some facts represented in the ontology; these will serve as examples to show how integrity checking of the information collected in the KB and the reasoning processes work.

Let us begin by considering the classes. As shown in the T-Box of Figure~\ref{fig:ont_animal}, the animals we can describe are: sardines, penguins, seals, polar bears, and reindeer. From T-Box we also see that we have only specified reindeer are herbivorous, and sardines are carnivorous, feeding habits of other spices will be inferred with hunt relationship.

We now examine some facts about the individuals that populate our KB:
\begin{itemize}
	\item Beary is a polar bear and hunt Sealy and Penguy;
	\item Sealy is a seal and hunt Penguy;
	\item Penguy is a penguin and hunt Sardy;
	\item Sardy is a sardine;
	\item Reiny is a reindeer and hunt Sardy; 
\end{itemize}

From this information, we can deduce that Beary and Reiny are predators; Sealy and Penguy are both predators and prey; and Sardy is prey.

\paragraph{Contraddiction: } A contradiction has also been deliberately introduced into the ontology. From the fact that Reiny hunts, it follows that she is a carnivorous animal; however, her species is labeled as herbivorous. Reiny is therefore both herbivorous and carnivorous. Moreover, herbivores and carnivores are disjoint classes. This makes the KB inconsistent.

If we used a reasoner to verify the consistency of the KB, for example HermiT\cite{glimm2014hermit}, we would obtain an error message informing us that the ontology is inconsistent and that the reasoner’s results cannot therefore be considered reliable.

The reasoner is also able to provide an explanation of the inconsistency (shown in Figure~\ref{fig:protege}\footnote{The figure shows the dialog window of Protégé\cite{musen2015protege}, a widely used ontology editor}), from which it is clear that the problem is that Reiny is simultaneously herbivorous and carnivorous.

\begin{figure}[h]
	\centering
	\includegraphics[width=0.9\linewidth]{quantum_ontology/protege}
	\caption{Error dialog in Protégé}
	\label{fig:protege}
\end{figure}

We introduced this inconsistency in order to show \textemdash{}in the following sections\textemdash{} how it is also possible in Prolog to detect this type of error and identify contradictions in the KB. If we remove the fact \say{Reiny hunt Sardy} the ontology will return consistent.

\section{Prolog version}



\section{Inference on the ontology}
\section{Conclusion}